\documentclass[conference]{IEEEtran}
\IEEEoverridecommandlockouts
% The preceding line is only needed to identify funding in the first footnote. If that is unneeded, please comment it out.
\usepackage{algorithm}
\usepackage{algpseudocode}
\usepackage{cite}
\usepackage{amsmath,amssymb,amsfonts}
%\usepackage{multirow}
\usepackage{graphicx}
\usepackage{textcomp}
\usepackage{xcolor}
\usepackage{color,soul}
\usepackage{float}
\usepackage{booktabs}
\usepackage{multirow}
\usepackage{subcaption}
\usepackage[acronym]{glossaries}
\usepackage{url} 
\usepackage[hidelinks]{hyperref}
\makeglossaries
%===========----------------------Glossaries----------------------===========
\loadglsentries{GL}
\renewcommand{\sectionautorefname}{Section}
\renewcommand{\subsectionautorefname}{Section}
\renewcommand{\tableautorefname}{Table}
\renewcommand{\figureautorefname}{Figure}
\def\BibTeX{{\rm B\kern-.05em{\sc i\kern-.025em b}\kern-.08em
    T\kern-.1667em\lower.7ex\hbox{E}\kern-.125emX}}
\begin{document}
\title{COAST: Coordinated Ownership and Anticipatory Sequential rouTing for Dynamic Vehicle Routing with Time Windows}
\author{
\IEEEauthorblockN{Tran Phuc}
\IEEEauthorblockA{\textit{Department of Computer Science, VNUHCM-University of Science}\\
Ho Chi Minh City, Vietnam\\
tdphuc.work@gmail.com}
}
	
\maketitle

\begin{abstract}
Dynamic Vehicle Routing Problems with Time Windows (DVRPTW) require sequential policies to simultaneously coordinate fleet-wide decisions and anticipate long-term routing consequences under uncertainty. Existing shared policies often entangle these two aspects within a single scoring mechanism, leading to suboptimal decisions and limited generalization. We propose COAST, a neural routing framework that explicitly decomposes customer selection into two signals: a coordination-driven ownership bias derived from per-vehicle memory, and a candidate-conditioned lookahead estimate of future cost. These signals are integrated through an edge-aware scoring function that captures dynamic feasibility. Experiments on DVRPTW benchmarks show that COAST outperforms classical heuristics and recent neural baselines. Ablation studies demonstrate that coordination and anticipation provide complementary benefits, while out-of-distribution evaluations confirm improved generalization across varying arrival rates, time-window constraints, and fleet sizes. Our results highlight the importance of explicitly separating coordination and anticipation in sequential routing policies for dynamic environments.
\end{abstract}

\begin{IEEEkeywords}
 Dynamic Vehicle Routing, Multi-Agent Reinforcement Learning, Structured Decision Decomposition, Coordination Memory, Anticipatory Planning
\end{IEEEkeywords}


\section{Introduction} \label{sec:intro}

The \gls{vrp} is a fundamental combinatorial optimization problem with wide applications in logistics and transportation systems. Classical \gls{vrp} formulations assume that all customer demands and travel information are known in advance. However, many real-world applications operate in dynamic environments where customer requests arrive over time, travel conditions may change, and routing decisions must be continuously updated. This leads to the \gls{dvrp}, an important extension of \gls{vrp} that better reflects real-world systems such as on-demand delivery and urban logistics.

In \gls{dvrptw}, each customer request is associated with temporal and capacity constraints, and decisions must be made online as new requests are revealed. Unlike static routing, where globally optimal solutions can be computed offline, dynamic settings require policies to make irrevocable decisions under partial information. Early decisions can significantly affect future system performance, creating a fundamental tension between immediate feasibility and long-term efficiency. This makes \gls{dvrptw} particularly challenging for traditional optimization methods that rely on complete and static problem information.

Existing approaches to \gls{dvrp} largely rely on rule-based heuristics, metaheuristics with periodic re-optimization, or rolling-horizon frameworks. While effective in specific settings, these methods often require substantial domain-specific tuning and lack robustness when applied to different dynamic regimes. Recently, deep reinforcement learning (DRL) has emerged as a promising paradigm for learning routing policies directly from data. By mapping system states to actions, DRL enables real-time decision making and has shown strong potential for combinatorial optimization problems.

Despite these advances, current learning-based approaches for \gls{dvrptw} still face fundamental limitations. A common paradigm is the shared sequential policy, where a single model makes decisions for all vehicles in turn. Although this enables implicit coordination through shared representations, it requires the policy to simultaneously address two distinct but tightly coupled challenges at each decision step: \textit{inter-vehicle coordination} (which vehicle should serve which customer) and \textit{anticipatory planning} (how current decisions affect future routing cost). 

Most existing architectures model both coordination and anticipation using a single compatibility score between vehicle and customer representations. This unified formulation forces the model to encode multiple decision factors within the same scoring function. As a result, the learning process suffers from three limitations. First, coordination and long-term evaluation impose competing learning signals on shared parameters, which can hinder stable optimization. Second, the learned policy often lacks robustness under distribution shifts, such as changes in request arrival patterns or time-window tightness. Third, the fused scoring mechanism reduces interpretability, making it difficult to distinguish whether suboptimal decisions arise from poor vehicle-customer assignment or inaccurate estimation of future routing cost.

We argue that these limitations stem from the lack of structural separation between coordination and anticipation. We hypothesize that explicitly decomposing these two components within the decision process can improve both routing performance and generalization. This is motivated by two observations: (i) in multi-agent systems, coordination can emerge from implicit communication via agent-specific memory, and (ii) in sequential decision-making, action-conditioned value estimation provides more stable guidance than global value functions in dynamic environments.

Based on this insight, we propose \gls{coast}, a neural routing framework that explicitly decomposes customer selection into two complementary signals. First, a coordination signal is derived from per-vehicle memory, where each vehicle maintains a latent state summarizing its routing history. An ownership predictor then produces soft assignment probabilities over customers, encouraging consistent fleet-level allocation. Second, an anticipation signal is modeled through a candidate-conditioned lookahead module that estimates the downstream cost of each feasible action. These signals are combined with a base compatibility score through an edge-aware fusion mechanism that incorporates dynamic feasibility constraints.

Importantly, our contribution lies not in introducing new architectural primitives, but in the principled decomposition of decision signals and the empirical demonstration of its effectiveness. We design hypothesis-driven experiments to isolate the roles of coordination and anticipation, and evaluate generalization under diverse dynamic settings.


The main contributions of this paper are summarized as follows:
\begin{itemize}
    \item We identify coordination-anticipation entanglement as a key limitation of existing sequential routing policies and propose a structured decomposition framework for \gls{dvrptw}.
    \item We introduce a memory-based coordination mechanism that enables implicit fleet allocation through learned ownership without explicit communication.
    \item We propose a candidate-conditioned lookahead module that captures long-term routing consequences at the action level.
    \item Extensive experiments demonstrate that the proposed approach improves routing performance and generalization compared to classical heuristics, hyper-heuristic and existing neural methods.
\end{itemize}

The remainder of this paper is organized as follows. \autoref{sec:related} reviews related work. \autoref{sec:prob} presents the problem formulation. \autoref{sec:method} describes the proposed DRL framework. \autoref{sec:experiment} reports computational results, and \autoref{sec:futurework} concludes the paper and outlines future research directions.
\section{Related works}\label{sec:related}

\section{Problem Formulation}
\label{sec:prob}

The Dynamic Vehicle Routing Problem with Time Windows (DVRPTW) extends the classical capacitated VRP with two additional sources of complexity: customer requests arrive dynamically during the execution of routes, and each request must be served within a prescribed time window.
Together, these extensions make it impossible to compute an optimal plan a priori and then execute it, since both the set of customers and the feasibility of serving them change as the fleet moves.
A routing policy must instead make real-time decisions for each vehicle --- selecting the next customer to serve based on the current state of the fleet and the currently known demand --- while balancing three competing criteria: minimizing the total distance traveled, minimizing accumulated tardiness from late arrivals, and maximizing the total number of customers served.
These objectives are in inherent tension: a policy that aggressively accepts marginal customers to maximize service rate may incur large detours and time-window violations; a policy that optimizes route efficiency may defer or skip customers it could have served with modest extra cost.
Managing this trade-off online, under incomplete information about future demand, constitutes the central challenge of DVRPTW.

We formalize the problem as a sequential Multi-agent Markov Decision Process (sMDDP)~\cite{TAM2026114234}, a variant of the multi-agent MDP framework that is adapted for settings where agents act asynchronously on a shared environment.
This model allows each vehicle to take decisions independently based on the global state of the fleet, while sharing a single learned policy that can generalize across problem instances without reoptimization.

\subsection{Problem Setting}
\label{sec:prob_setting}

\subsubsection{Instance structure}
The instance is defined over a node set $\{0, 1, \ldots, L_c\}$.
Node $j = 0$ denotes the \emph{depot}: the common origin and termination point for all vehicle routes.
The depot has no demand, no service duration, and no time-window constraint; it is not counted as a served customer in the objective.
Nodes $j \geq 1$ represent customer requests, each associated with a geographic location, a demand quantity, a service time, and a time window.

\subsubsection{Customer features}
Each customer $j \geq 1$ is described by a static feature vector
\[
\mathbf{x}_j = \bigl(x_j,\; y_j,\; d(j),\; tw^{start}_j,\; tw^{end}_j,\; s(j)\bigr) \in \mathbb{R}^{D_c},
\]
where $(x_j, y_j)$ is the normalized geographic location, $d(j)$ the demand (normalized by vehicle capacity), $[tw^{start}_j, tw^{end}_j]$ the time window within which service must \emph{begin}, and $s(j)$ the service duration.
The time window $[tw^{start}_j, tw^{end}_j]$ encodes the customer's schedule constraint: if a vehicle arrives before $tw^{start}_j$, it must wait; if it arrives after $tw^{end}_j$, the delivery is late.
In this work we adopt \emph{soft} time windows following~\cite{TAM2026114234}: a vehicle may arrive after $tw^{end}_j$, but any excess is penalized proportionally in the objective.
This relaxation is motivated by practical logistics scenarios in which a late delivery incurs a cost penalty rather than an absolute rejection, and it also enables the policy to learn a smoother reward landscape during training compared to hard infeasibility constraints.

\subsubsection{Vehicle state}
At any decision step $t$, vehicle $k$ is described by a continuous state vector
\[
\mathbf{s}_k^t = \bigl(x_k,\; y_k,\; q^{rem}_k,\; t^{cur}_k\bigr) \in \mathbb{R}^4.
\]
Here $(x_k, y_k)$ is the vehicle's current position (the location of the most recently served customer or the depot), $q^{rem}_k$ is the remaining cargo capacity after all committed deliveries, and $t^{cur}_k$ is the vehicle's \emph{next availability time} --- the earliest time at which it can depart toward its next customer, accounting for the travel time to reach its current position plus any service time that has not yet elapsed.
This compact four-dimensional state is sufficient to determine the vehicle's feasibility with respect to any pending customer: it encodes where the vehicle is, how much it can carry, and when it will be free.

\subsubsection{Travel, waiting, arrival, and tardiness}
When vehicle $k$ decides to serve customer $j$, four quantities determine the cost and feasibility of the assignment:
\begin{align}
d_{kj}   &= \bigl\|(x_k, y_k) - (x_j, y_j)\bigr\|_2,                       \label{eq:dist}\\
\tau_{kj} &= d_{kj} / v,                                                      \label{eq:ttravel}\\
a_{kj}   &= t^{cur}_k + \tau_{kj},                                            \label{eq:arrival}\\
\ell_{kj} &= \max\!\bigl(0,\; a_{kj} - tw^{end}_j\bigr).                     \label{eq:tardiness}
\end{align}
Equation~\eqref{eq:dist} gives the Euclidean travel distance; equation~\eqref{eq:ttravel} converts it to travel time at constant speed $v$; equation~\eqref{eq:arrival} gives the arrival time at $j$; and equation~\eqref{eq:tardiness} defines the \emph{tardiness} $\ell_{kj} \geq 0$ as the excess beyond the closing time.
If $a_{kj} < tw^{start}_j$, the vehicle must wait $tw^{start}_j - a_{kj}$ time units before service can begin; this waiting time is absorbed into $t^{cur}_{k}$ after the assignment.
A customer $j$ is considered \emph{capacity-feasible} for vehicle $k$ if $d(j) \leq q^{rem}_k$.
Under soft time windows, every revealed customer with $d(j) \leq q^{rem}_k$ is a valid candidate action regardless of tardiness; the magnitude of $\ell_{kj}$ directly enters the reward, so the policy learns to avoid late arrivals without being constrained by hard cutoffs.

\subsubsection{Dynamic arrivals}
A distinguishing feature of DVRPTW is that the full customer set is not known at the start of the planning horizon.
Customer requests are partitioned into an initial visible set $\mathcal{C}_0$ known at $t=0$, and latent batches $\mathcal{C}_{e_1}, \mathcal{C}_{e_2}, \ldots$ that are revealed at event epochs $e_1 < e_2 < \cdots$ during route execution.
The proportion of customers that are initially hidden is characterized by the \emph{degree of dynamism} (DoD): $\mathrm{DoD} = 0$ corresponds to the fully static VRPTW (all customers known a priori), while $\mathrm{DoD} = 1$ means every customer arrives dynamically and no advance planning is possible.
At each decision step, the policy can observe only the currently revealed and unserved customers; future arrivals are unknown.
This partial observability forces the policy to balance exploiting the current known demand against preserving enough temporal and capacity slack to accommodate customers that may appear later.
A high DoD requires more adaptive routing strategies, as vehicles cannot commit to long fixed routes without risking infeasibility upon new arrivals.

\subsection{Sequential Multi-Agent Decision Model}
\label{sec:prob_mdp}

\subsubsection{Motivation for the sMDDP framework}
Modeling DVRPTW as a standard single-agent MDP with a centralized policy would require producing joint actions for all vehicles simultaneously at each global time step.
This is problematic for two reasons.
First, the joint action space $\prod_{k=1}^{L_v} \mathcal{A}^k$ grows exponentially with fleet size, making direct optimization intractable beyond small fleets.
Second, actions for different vehicles have different durations: a vehicle that commits to a nearby customer becomes available again long before a vehicle that travels across the map.
Forcing all vehicles to act in a synchronized step ignores this temporal structure and would require artificial padding or discretization of time.

An alternative is to treat each vehicle as a fully independent agent with its own policy, ignoring the other vehicles.
However, this fails to capture the essential coordination problem: the optimal route for vehicle $k$ depends on what routes the other vehicles are committed to, since they are competing for the same customers.
Without observing the full fleet state, individual policies cannot avoid duplication or leave regions underserved.

The sMDDP~\cite{TAM2026114234} provides a principled middle ground.
It models the multi-vehicle problem as a sequence of individual decisions, each made by one vehicle at the moment it becomes available, conditioned on the full global state of the fleet.
This preserves the asynchronous, variable-duration nature of vehicle actions while making each decision tractable, and it allows a single shared policy to coordinate implicitly through the global state.

\subsubsection{Formal definition}
We define the DVRPTW sMDDP as the tuple $(\mathcal{I}, \mathcal{S}, \mathcal{A}, P, P_0, R, T)$:

\smallskip\noindent\textbf{Agents.}
$\mathcal{I} = \{1, \ldots, L_v\}$ is the set of vehicle agents.
All agents share the same action structure and reward signal, enabling a homogeneous policy.

\smallskip\noindent\textbf{State space.}
The global state $s_t \in \mathcal{S}$ is factored as
\[
s_t = \bigl(\{\mathbf{s}_k^t\}_{k \in \mathcal{I}},\;\; s_t^E,\;\; \iota_t\bigr),
\]
where $\{\mathbf{s}_k^t\}$ collects the individual vehicle states, $s_t^E$ is the \emph{joint environment state} encoding the feature vector $\mathbf{x}_j$ and the current service status (hidden, pending, or served) of every customer that has been revealed so far, and $\iota_t \in \mathcal{I}$ is the \emph{turn variable} identifying which vehicle must make a decision at step $t$.
The factored structure separates information that is \emph{local} to each vehicle (position, capacity, availability time) from information that is \emph{shared} across the fleet (customer status), enabling the policy to exploit both levels of context when computing a routing decision.

\smallskip\noindent\textbf{Action space.}
At step $t$, the active vehicle $\iota_t$ selects
\[
a_t \in \mathcal{A}(s_t, \iota_t),
\]
where $\mathcal{A}(s_t, \iota_t)$ is the \emph{individual feasible action set}: the set of customers that are currently pending (revealed and unserved) and satisfy the capacity constraint $d(j) \leq q^{rem}_{\iota_t}$.
Returning to the depot (action $a_t = 0$) is always available and marks the end of vehicle $\iota_t$'s route for that episode.

\smallskip\noindent\textbf{Transition function.}
$P : \mathcal{S} \times \mathcal{A} \times \mathcal{S} \to [0,1]$ governs state evolution.
Executing action $a_t$ triggers four updates: (i) vehicle $\iota_t$'s position moves to $(x_{a_t}, y_{a_t})$ and its next availability time advances to $a_{\iota_t, a_t} + s(a_t)$; (ii) capacity decreases by $d(a_t)$; (iii) customer $a_t$ is marked as served in $s_t^E$; and (iv) any customer batches whose reveal epoch falls before the new time $t^{cur}_{\iota_t}$ become visible and are added to $s_t^E$ as pending.
The next active vehicle $\iota_{t+1}$ is the vehicle with the smallest next availability time among all vehicles that still have pending customers or have not yet returned to the depot.

\smallskip\noindent\textbf{Initial state and episode length.}
$P_0$ defines the initial distribution: all vehicles begin at the depot at time $0$ with full capacity.
The episode length $T$ is variable and ends when all vehicles have returned to the depot.

\subsubsection{Asynchronous sequential structure}
The turn variable $\iota_t$ is the mechanism that converts the asynchronous multi-vehicle problem into a well-defined sequential process.
Because traveling to and serving a customer takes a duration proportional to distance plus service time, different vehicles become available at different moments.
Rather than advancing a global clock and forcing all vehicles to act simultaneously, the sMDDP advances the decision index $t$ each time any one vehicle becomes ready.
The ordering of decisions is therefore determined by the data: a vehicle that picks a distant customer will wait longer before acting again than one that picks a nearby customer.

This event-driven structure has an important implication for policy design: at the moment vehicle $\iota_t$ acts, the positions and availability times of all other vehicles are known exactly (they are part of $s_t$), even though those vehicles have not yet reached their next destinations.
The policy therefore has access to a synchronized snapshot of the full fleet's committed state, which is the information needed to avoid routing conflicts, such as two vehicles targeting the same customer simultaneously.

\subsubsection{Shared homogeneous policy}
Since all vehicles perform the same routing task under the same reward signal, we adopt a homogeneous parameterization: all vehicles share a single policy network $\pi_\theta(a_t \mid s_t, \iota_t)$.
At each step, the policy receives the full global state $s_t$ and the identity of the active vehicle $\iota_t$, and outputs a probability distribution over $\mathcal{A}(s_t, \iota_t)$.
Weight sharing across agents reduces the parameter count by a factor of $L_v$ compared to independent per-vehicle networks, and it endows the policy with a form of inductive bias: strategies learned by one vehicle in one context are automatically transferred to all other vehicles in analogous contexts.
Critically, this parameterization supports any fleet size $L_v$ at inference time without retraining, since the policy maps a fixed-size state representation to a score for each individual candidate --- the fleet size affects only the number of candidates scored per step, not the network architecture.
As shown in~\cite{TAM2026114234}, the separability of the transition and reward functions in this model ensures that optimizing the shared policy is equivalent to solving the joint optimization problem, provided the fleet is homogeneous.

\subsection{Objective}
\label{sec:prob_obj}

\subsubsection{Multi-criteria cost structure}
The DVRPTW routing objective incorporates three distinct and competing cost components.

\textit{Travel distance.}
Every arc traversed by any vehicle incurs a routing cost proportional to its Euclidean length.
Minimizing total distance is the classical VRP objective and captures the primary operational cost of deploying a fleet.

\textit{Tardiness.}
When a vehicle arrives at a customer after its closing time $tw^{end}_j$, the excess delay $\ell_{kj} = \max(0, a_{kj} - tw^{end}_j)$ represents a service-quality failure.
Accumulating tardiness across all deliveries constitutes the second cost component.
Reducing tardiness often conflicts with minimizing distance: respecting tight time windows may require a vehicle to travel to a customer earlier in the route even if it is geographically less efficient to do so, or to prioritize a time-critical customer over a closer but less urgent one.

\textit{Service acceptance.}
A vehicle's route must be feasible with respect to remaining capacity; customers whose demand would exceed the residual cargo cannot be served and must be rejected.
In the dynamic setting, rejecting a customer is irreversible: once the vehicle has passed the window in which the detour would have been feasible, the customer cannot be recovered.
Maximizing the number of served customers is therefore a separate objective that may conflict with both distance and tardiness minimization, as accepting a marginal customer may require a costly detour or a time-window violation.

\subsubsection{Scalarized reward formulation}
To enable training with a single policy via reinforcement learning, we scalarize the three objectives into a per-step reward.
At each step $t$ in which vehicle $\iota_t$ serves customer $a_t$, the immediate reward is
\[
r_t = -\,d_{\iota_t, a_t} \;-\; \lambda_{\mathrm{late}}\,\ell_{\iota_t, a_t},
\]
capturing the travel cost and tardiness incurred by the assignment.
At the end of the episode, an additional terminal penalty accounts for unserved demand:
\[
r_T^{\mathrm{skip}} = -\,\lambda_{\mathrm{skip}}\, N_{\mathrm{skip}},
\]
where $N_{\mathrm{skip}}$ counts all revealed but unserved customer requests.
The total undiscounted return for one episode is therefore:
\begin{align*}
R &= \sum_{t=0}^{T-1} r_t + r_T^{\mathrm{skip}} \\
  &= -\sum_t d_{\iota_t, a_t}
  \;-\; \lambda_{\mathrm{late}} \sum_t \ell_{\iota_t, a_t}
  \;-\; \lambda_{\mathrm{skip}}\, N_{\mathrm{skip}}.
\end{align*}
This return corresponds to the negative of the evaluation cost $C$ used in Section~\ref{sec:experiment}, ensuring that training and evaluation are optimizing the same quantity.
The hyperparameters $\lambda_{\mathrm{late}}$ and $\lambda_{\mathrm{skip}}$ control the relative priority of time-window compliance and service completeness versus routing efficiency.
Large $\lambda_{\mathrm{late}}$ produces a policy that prioritizes on-time delivery even at the cost of longer routes; large $\lambda_{\mathrm{skip}}$ produces a policy that accepts customers aggressively; small values of both hyperparameters recover a distance-only routing policy similar to classical CVRP.

The training objective is to find parameters $\theta$ that maximize the expected return over the training distribution:
\[
\max_\theta \;\mathbb{E}_{\pi_\theta}\!\left[R\right].
\]

\subsubsection{Key challenges}
This formulation surfaces three structural properties of DVRPTW that must be addressed by any effective policy.
\textit{Non-stationarity}: the set of actionable customers changes at every step due to new arrivals, vehicle movements, and constraint changes, making it impossible to compute a fixed optimal assignment in advance.
\textit{Sequential coupling}: each vehicle's decision immediately affects the feasible options available to the remaining vehicles, introducing an implicit coordination requirement that grows with fleet size and route length.
\textit{Interaction-dependent feasibility}: both the tardiness $\ell_{kj}$ and the capacity gap depend on the specific state of vehicle $k$ at the moment it would serve customer $j$ --- they cannot be precomputed from static node features alone, but must be evaluated dynamically for each vehicle--customer pair at each decision step.
These three properties collectively motivate the edge-aware compatibility scoring, memory-based fleet coordination, and candidate-conditioned lookahead mechanisms introduced in Section~\ref{sec:method}.
% !TEX root = main_AAAI.tex
\section{Methodology}\label{sec:method}

This section presents the COAST architecture and its training procedure.
Table~\ref{tab:notation} summarizes the main notation.

\subsection{Policy Overview}\label{sec:policy_overview}

Existing neural routing policies typically produce a selection score through a single compatibility function between vehicle and customer representations~\cite{kool2019attention}, which conflates immediate feasibility assessment with fleet-level coordination and future planning.
When these signals interfere, the policy gradient receives mixed learning signals that slow convergence and limit generalization under distribution shifts.

COAST instead decomposes the customer-selection score into three independently computed signals combined by a learned fusion mechanism.
At each step $t$, given state $s_t$ and active vehicle $k_t$, the policy is:
\[
\pi_\theta(a_t \mid s_t, k_t)
=
\frac{\exp\!\big(S_\theta(a_t \mid s_t, k_t)\big)}
{\displaystyle\sum_{u \in \mathcal{A}(s_t, k_t)} \exp\!\big(S_\theta(u \mid s_t, k_t)\big)},
\]
where $\mathcal{A}(s_t, k_t)$ is the feasible action set and the composite score is:
\[
S_\theta(j \mid s_t, k_t) = \phi_{fuse}\!\left(\tilde{S}^{att}_j,\; \tilde{S}^{owner}_j,\; \tilde{S}^{look}_j\right).
\]
The three signals are:
\begin{itemize}
    \item $S^{att}$: a \textit{local compatibility score} quantifying vehicle--customer fit given current spatial, temporal, and capacity constraints;
    \item $S^{owner}$: a \textit{coordination bias} derived from per-vehicle memory, encoding implicit fleet-level assignment of customer $j$ to vehicle $k_t$;
    \item $S^{look}$: a \textit{candidate-conditioned lookahead estimate} of the downstream routing consequence of committing to customer $j$.
\end{itemize}

\subsection{Model Architecture}\label{sec:model_architecture}

\paragraph{Customer graph encoding.}
We encode customers using a Transformer-based graph encoder~\cite{vaswani2017attention} that produces contextual embeddings capturing inter-customer relational structure.
Depot and customer nodes use separate linear projections since they play structurally different roles:
\[
\mathbf{h}_j^{(0)} =
\begin{cases}
\mathbf{W}_{dep}\,\mathbf{x}_j, & j = 0 \\
\mathbf{W}_{cust}\,\mathbf{x}_j, & j \geq 1
\end{cases}
\]
where $\mathbf{x}_j \in \mathbb{R}^{D_c}$ is the raw feature vector.
Since geographic proximity is a strong signal for routing quality but may not emerge reliably from content-based attention alone, we augment each attention layer with a distance-aware bias: a 16-bin RBF encoding of normalized pairwise distance $d_{ij} \in [0,1]$,
\[
\mathrm{RBF}_k(d) = \exp\!\left(-\frac{(d - \mu_k)^2}{2\sigma^2}\right),
\]
with 16 evenly-spaced centres $\mu_k = (k-1)/15$, $k=1,\ldots,16$, and bandwidth $\sigma = 1/15$.
The 16-dimensional encoding is projected by an MLP to a per-head scalar bias $b_{ij}^{(l,h)}$ added to the attention logit.
Each layer applies two post-norm residual sublayers~\cite{vaswani2017attention}.
Letting $\mathrm{MHA}_i^{(l)} = \mathbf{W}_{out}^{(l)}\!\left[\bigoplus_{h=1}^{H}\sum_j \alpha_{ij}^{(l,h)}\mathbf{v}_j^{(l,h)}\right]$ denote the multi-head attention output:
\begin{align*}
\hat{\mathbf{h}}_i^{(l)} &= \mathrm{LN}\!\left(\mathbf{h}_i^{(l)} + \mathrm{MHA}_i^{(l)}\right), \\
\mathbf{h}_i^{(l+1)} &= \mathrm{LN}\!\left(\hat{\mathbf{h}}_i^{(l)} + \mathrm{FFN}\!\left(\hat{\mathbf{h}}_i^{(l)}\right)\right).
\end{align*}
The final customer representation is $\mathbf{c}_j = \mathrm{Dropout}(\mathbf{W}_c\,\mathbf{h}_j^{(L)}) \in \mathbb{R}^D$.
In the dynamic setting, encodings are recomputed whenever new orders become visible.

\paragraph{Acting vehicle encoding.}
Only the active vehicle $k_t$ is encoded at each step --- full fleet re-encoding at every step is wasteful and unnecessary since inter-vehicle coordination is handled by CoordinationMemory rather than vehicle-to-vehicle attention.
The vehicle state $\mathbf{s}_{k_t} = (x, y, q^{rem}, t^{cur}) \in \mathbb{R}^4$ is projected and refined through $L$ layers of cross-attention over $H^{cust}$, restricted to the feasible action set:
\[
h^{act} = \mathbf{W}_v\;\mathrm{FleetEncoder}\!\bigl(\mathbf{W}_{in}\,\mathbf{s}_{k_t},\; H^{cust},\; \mathcal{A}(s_t,k_t)\bigr),
\]
where $h^{act} \in \mathbb{R}^{N \times 1 \times D}$.

\paragraph{Edge feature encoding}
Node-level representations capture global structure but cannot express the state-dependent interaction between a specific vehicle and a specific candidate at this step.
Two customers equidistant from the vehicle may differ radically in time-window compatibility, waiting time, or remaining capacity fit.
We therefore construct an 8-dimensional edge feature vector $\mathbf{e}_{ij} \in \mathbb{R}^8$ for each (vehicle $k_t$, candidate $j$) pair:

\medskip
\begin{center}
\begin{tabular}{@{} l l p{1.65in} @{}}
\toprule
Feature & Symbol & Definition \\
\midrule
Distance       & $d_{ij}$    & $\|\mathbf{p}_{k_t} - \mathbf{p}_j\|_2$ \\
Travel time    & $\tau_{ij}$ & $d_{ij} / v$ \\
Arrival time   & $a_{ij}$    & $t^{cur}_{k_t} + \tau_{ij}$ \\
Waiting time   & $w_{ij}$    & $\max(tw^{start}_j - a_{ij},\,0)$ \\
Lateness       & $l_{ij}$    & $\max(a_{ij} - tw^{end}_j,\,0)$ \\
Time slack     & $s_{ij}$    & $tw^{end}_j - a_{ij}$ \\
Feasibility    & $f_{ij}$    & 1 if $j$ is capacity- and time-feasible, else 0 \\
Capacity gap   & $g_{ij}$    & $q^{rem}_{k_t} - d(j)$ \\
\bottomrule
\end{tabular}
\end{center}
\medskip

Together these features encode whether the vehicle arrives on time, how much schedule slack remains, and whether capacity is sufficient --- information that changes at every step as the vehicle's state evolves.
The features are encoded through a two-layer MLP with layer normalization:
\[
\mathbf{e}_{ij}^{emb} = \mathrm{LN}\!\left(\mathbf{W}_{e2}\,\mathrm{ReLU}\!\left(\mathbf{W}_{e1}\,\mathbf{e}_{ij}\right)\right) \in \mathbb{R}^D,
\]
yielding the full edge tensor $H^{edge} \in \mathbb{R}^{N \times 1 \times L_c \times D}$.

\paragraph{Local compatibility via edge-aware scoring.}
Standard attention scoring~\cite{kool2019attention} computes vehicle--customer compatibility from node-level representations that are oblivious to the specific interaction at this step.
The CrossEdgeFusion module injects the per-pair edge embedding as an additive bias into multi-head attention logits:
\[
u_j^{(h)} = \frac{\mathbf{q}^{(h)} \cdot \mathbf{k}_j^{(h)}}{\sqrt{d_h}} + \mathbf{e}_j^{emb} \cdot \mathbf{w}_{bias}^{(h)},
\]
where $\mathbf{q}^{(h)} = h^{act}\mathbf{W}_Q^{(h)}$ and $\mathbf{k}_j^{(h)} = \mathbf{c}_j\mathbf{W}_K^{(h)}$.
The compatibility score averages the pre-softmax logits across heads:
\[
S^{att}_j = \frac{1}{H}\sum_{h=1}^{H} u_j^{(h)} \in \mathbb{R}^{N \times 1 \times L_c}.
\]
No softmax is applied inside CrossEdgeFusion; the single softmax over all candidates is deferred to the final action selection stage after all three signals are fused, ensuring that inter-candidate probability mass is distributed over the most informative composite signal.

\paragraph{Coordination through memory and ownership.}
Without a coordination mechanism, the policy may route multiple vehicles toward the same customer clusters while neglecting others.
Explicit fleet-level communication at every step is computationally heavy and unnecessary when softer implicit signals suffice.

We instead equip each vehicle with a private latent memory $\mathbf{m}_k^{(t)} \in \mathbb{R}^{H_m}$, initialized to zero, that summarizes its routing history.
After vehicle $k_t$ selects customer $j^*$, its memory is updated via a bounded, GRU-like recurrence:
\[
\mathbf{m}_{k_t}^{(t+1)} = \tanh\!\left(\mathbf{W}_{in}\!\left[\,h^{act},\;\mathbf{c}_{j^*},\;\mathbf{e}_{k_t j^*}^{emb}\,\right] + \mathbf{W}_{hid}\,\mathbf{m}_{k_t}^{(t)}\right).
\]
Only the acting vehicle's memory slot is updated; all others remain unchanged.

This per-vehicle memory supports an OwnershipHead that computes a soft fleet-level customer assignment.
A vehicle whose memory is strongly aligned with customer $j$ has implicitly claimed it through prior routing choices.
The assignment logits are computed via a bilinear interaction over all vehicles:
\[
o_{ik} = \frac{(\mathbf{W}_{veh}\,\mathbf{m}_i^{(t)}) \cdot (\mathbf{W}_{cust}\,\mathbf{c}_k)^\top}{\sqrt{D}},
\]
and the coordination bias for vehicle $k_t$ is the log-relative ownership over the fleet:
\[
S^{owner}_j
= \log\!\left(\frac{\exp(o_{k_t,j})}{\sum_{i'=1}^{L_v}\exp(o_{i',j})}\right),
\quad S^{owner} \in \mathbb{R}^{N \times 1 \times L_c}.
\]
Using log-probability ensures the ownership signal is well-scaled relative to the other scores, and the softmax across vehicles encodes \textit{relative} claim --- a high $S^{owner}_j$ means vehicle $k_t$ is the strongest claimant among all vehicles.
Neither the OwnershipHead nor the memory update imposes any auxiliary supervised loss; both are shaped entirely by the policy gradient, allowing coordination semantics to emerge from optimizing routing performance.

\paragraph{Candidate-conditioned lookahead estimation.}
A policy that maximizes immediate feasibility may still be myopic: selecting a nearby customer can leave the vehicle in a position that violates time windows of unserved neighbors, or forces an early depot return due to capacity exhaustion.
A global state value $V(s_t)$ cannot differentiate between candidates and therefore offers no guidance at the action level.

The LookaheadHead produces a per-candidate scalar that estimates the downstream routing consequence of committing to customer $j$:
\[
S^{look}_j = \mathbf{W}_{\ell 2}\;\mathrm{Dropout}\!\bigl(\mathrm{ReLU}\bigl(\mathbf{W}_{\ell 1}
\bigl[\,h^{act},\;\mathbf{c}_j,\;\mathbf{e}_{ij}^{emb}\,\bigr]\bigr)\bigr),
\]
producing a per-candidate scalar ($S^{look}_j \in \mathbb{R}$) conditioned on the vehicle's current state, the candidate's global embedding, and their dynamic interaction.
Critically, this head carries no explicit regression target --- $S^{look}$ is shaped entirely through the policy gradient, which drives the fusion MLP to weight lookahead appropriately when it reduces long-term routing cost.

\paragraph{Multi-signal score fusion.}
The three signals are computed by different mechanisms and operate on incompatible numerical scales.
To normalize them before fusion, we apply mask-aware z-normalization over the feasible candidates for each signal independently.
For each $j \in \mathcal{A}(s_t, k_t)$:
\[
\tilde{S}^{(\cdot)}_j = \frac{S^{(\cdot)}_j - \mu^{(\cdot)}}{\sqrt{(\sigma^{(\cdot)})^2 + \varepsilon}},
\]
where $\mu^{(\cdot)}$ and $\sigma^{(\cdot)}$ are the mean and standard deviation computed \emph{only} over feasible candidates; scores for infeasible positions are subsequently set to $-\infty$ before the final softmax.

The normalized signals are fused by a compact two-layer MLP:
\[
S^{final}_j = \mathbf{W}_{f2}\;\mathrm{ReLU}\!\left(\mathbf{W}_{f1}\!\left[\,\tilde{S}^{att}_j,\;\tilde{S}^{owner}_j,\;\tilde{S}^{look}_j\,\right]\right),
\]
with $\mathbf{W}_{f1} \in \mathbb{R}^{64 \times 3}$ and $\mathbf{W}_{f2} \in \mathbb{R}^{1 \times 64}$.
A fixed-weight sum cannot adapt to the fact that ownership is uninformative early in an episode but becomes the dominant coordination cue later; the MLP learns these context-dependent interaction weights through the policy gradient.
An optional tanh scaling bounds the magnitude: $S^{final}_j \leftarrow \beta_0 \cdot \tanh(S^{final}_j)$.

\subsection{Training Objective}\label{sec:training_objective}

COAST is trained end-to-end using REINFORCE~\cite{nazari2018reinforcement} with a learned critic baseline.
The critic $b_\phi$ and the LookaheadHead play strictly separate roles: $b_\phi$ is a state-level variance-reducing baseline that is agnostic to the action chosen, while $S^{look}$ is an action-conditioned routing guide.
Conflating them would couple value estimation with the scoring mechanism, degrading both.

\paragraph{Policy gradient with critic baseline.}
The policy loss is:
\[
\mathcal{L}_{policy}(\theta)
=
-\mathbb{E}_{\tau \sim \pi_\theta}\!\left[\sum_{t=0}^{T} \log \pi_\theta(a_t \mid s_t, k_t)\, A_t \right],
\]
where $A_t = R_t - b_\phi(s_t)$ is the advantage, $R_t = \sum_{k=t}^T r_k$ is the cumulative return, and the critic is trained with:
\[
\mathcal{L}_{critic}(\phi)
=
\mathbb{E}\!\left[\sum_{t=0}^{T} \mathrm{SmoothL1}\!\left(b_\phi(s_t),\; R_t\right)\right].
\]
Smooth-L1 is used for robustness to the return variance induced by stochastic customer arrivals.
Advantage estimates are batch-normalized before computing $\mathcal{L}_{policy}$.

\paragraph{Combined objective.}
We augment with an entropy term to prevent premature convergence across the diverse instance distribution:
\begin{align*}
\mathcal{L}(\theta, \phi)
&= \mathcal{L}_{policy}(\theta)
+ \alpha\; \mathcal{L}_{critic}(\phi) \\
&\quad -\,\beta\;\mathbb{E}\!\left[\sum_{t=0}^{T}
  \mathcal{H}\!\left(\pi_\theta(\cdot \mid s_t, k_t)\right)\right],
\end{align*}
with $\beta = 0.01$.
The LookaheadHead and OwnershipHead carry no auxiliary losses; both receive gradient signal exclusively through $S^{final}$ via the policy gradient.

\subsection{Discussion}\label{sec:method_discussion}

The design of COAST reflects three structural inductive biases aligned with dynamic multi-vehicle routing.

\textit{Edge features as first-class routing context.}
Static node embeddings encode global relational structure but cannot express how vehicle state and customer constraints interact at a given step.
The explicit construction of per-pair edge features --- encoding time slack, lateness, capacity gap, and feasibility --- allows the scoring mechanism to react immediately to changes in vehicle state without waiting for global re-encoding.
This is the primary mechanism through which COAST achieves constraint-awareness at decision time.

\textit{Communication-free fleet coordination via emergent ownership.}
Explicit inter-vehicle communication or full fleet re-encoding at every step adds overhead proportional to fleet size and introduces gradient coupling between vehicles.
The per-vehicle memory and OwnershipHead provide a scalable alternative: coordination emerges implicitly from each vehicle's trajectory history, naturally inducing spatial specialization without any vehicle-to-vehicle attention.
The ownership signal is shaped purely by routing performance, not by any auxiliary supervision, so the coordination structure it encodes reflects what actually improves fleet efficiency.

\textit{Action-conditioned lookahead without a supervised target.}
The LookaheadHead differs from standard value-based approaches in that it produces per-candidate estimates without being constrained to approximate any specific quantity such as $Q(s_t, a_t)$.
By training entirely through the policy gradient, the lookahead signal is free to encode whatever combination of vehicle state, candidate embedding, and edge interaction best reduces myopic decisions --- including routing consequences that are difficult to formalize analytically.
Its separation from the critic baseline preserves the action-independence required for an unbiased gradient estimator while allowing the per-candidate signal to carry independent predictive content.

\begin{table}[t]
\centering
\caption{Summary of main notations used in Section~\ref{sec:method}.}
\label{tab:notation}
\resizebox{\linewidth}{!}{%
\begin{tabular}{ll}
\toprule
\textbf{Symbol} & \textbf{Description} \\
\midrule

\multicolumn{2}{l}{\textit{Decision process}} \\
$t$                          & Decision step index \\
$s_t$                        & Environment state at step $t$ \\
$k_t$                        & Index of the active vehicle at step $t$ \\
$a_t$                        & Selected customer (action) \\
$\mathcal{A}(s_t, k_t)$      & Feasible action set for vehicle $k_t$ \\

\midrule
\multicolumn{2}{l}{\textit{Problem dimensions}} \\
$N$   & Batch size \\
$L_c$ & Number of customers including depot \\
$L_v$ & Number of vehicles \\
$D$   & Embedding dimension \\
$H$   & Number of attention heads \\
$H_m$ & Coordination memory hidden size \\
$H_\ell$ & Lookahead MLP hidden size \\

\midrule
\multicolumn{2}{l}{\textit{Input features}} \\
$\mathbf{x}_j \in \mathbb{R}^{D_c}$      & Raw feature vector of customer or depot $j$ \\
$\mathbf{s}_{k_t} \in \mathbb{R}^4$       & State of active vehicle: $(x, y, q^{rem}, t^{cur})$ \\
$\mathbf{e}_{ij} \in \mathbb{R}^8$        & Raw edge feature vector for pair (vehicle $k_t$, candidate $j$) \\

\midrule
\multicolumn{2}{l}{\textit{Learned representations}} \\
$\mathbf{c}_j \in \mathbb{R}^D$,\; $H^{cust} \in \mathbb{R}^{N \times L_c \times D}$               & Customer embeddings from GraphEncoder \\
$h^{act} \in \mathbb{R}^{N \times 1 \times D}$                                                       & Active vehicle embedding (one vehicle encoded per step) \\
$\mathbf{e}_{ij}^{emb} \in \mathbb{R}^D$,\; $H^{edge} \in \mathbb{R}^{N \times 1 \times L_c \times D}$ & Encoded edge features \\

\midrule
\multicolumn{2}{l}{\textit{Scoring signals}} \\
$S^{att}   \in \mathbb{R}^{N \times 1 \times L_c}$ & Local compatibility score (CrossEdgeFusion) \\
$S^{owner} \in \mathbb{R}^{N \times 1 \times L_c}$ & Ownership coordination bias (log-probability) \\
$S^{look}  \in \mathbb{R}^{N \times 1 \times L_c}$ & Candidate-conditioned lookahead estimate \\
$S^{final} \in \mathbb{R}^{N \times 1 \times L_c}$ & Fused routing score \\
$\tilde{S}^{(\cdot)}$                               & Z-normalized version of each signal \\

\midrule
\multicolumn{2}{l}{\textit{Coordination memory}} \\
$\mathbf{m}_k^{(t)} \in \mathbb{R}^{H_m}$    & Memory state of vehicle $k$ at step $t$ \\
$M_t \in \mathbb{R}^{N \times L_v \times H_m}$ & Full fleet memory tensor at step $t$ \\

\midrule
\multicolumn{2}{l}{\textit{Policy and learning}} \\
$\pi_\theta(a_t \mid s_t, k_t)$  & Routing policy parameterized by $\theta$ \\
$R_t$                            & Cumulative return from step $t$ \\
$b_\phi(s_t) \in \mathbb{R}$     & Learned critic baseline (distinct from LookaheadHead) \\
$A_t = R_t - b_\phi(s_t)$        & Advantage estimate \\
$\alpha,\,\beta$                 & Critic weight and entropy coefficient \\
\bottomrule
\end{tabular}}
\end{table}

\section{Experimental results} \label{sec:experiment}

We evaluate COAST on two complementary benchmark families to assess routing quality, dynamic robustness, and the contribution of each architectural component. The primary metric is normalized routing cost under greedy decoding; sampling-based decoding is reported where available as a supplementary reference. Unless otherwise stated, lower cost is better.

\subsection{Experimental setup}

Our experimental study is organized around three complementary questions. First, we evaluate whether COAST improves solution quality relative to neural and classical baselines on established DVRPTW benchmark families. Second, we examine robustness under varying degrees of dynamism and time-window tightness to assess whether the proposed decomposition remains effective beyond a single operating regime. Third, we perform targeted ablations to isolate the contributions of ownership, lookahead, edge-aware scoring, and multi-signal fusion. The remainder of this subsection therefore introduces the evaluation metric, benchmark instances, baselines, and training configuration that support these analyses.

\paragraph{Evaluation metric}
For each test instance, we report the total normalized routing cost of the decoded solution. Let $d_t$ denote the travel distance incurred by the acting vehicle at decision step $t$, and let $s_t^{\mathrm{start}}$ be the service start time at the selected customer. If the selected customer has closing time $b_t$, then the stepwise lateness is defined as
\[
\ell_t = \max(0, s_t^{\mathrm{start}} - b_t).
\]
The cost of one instance is
\[
C
=
\sum_t d_t
\;+\;
\lambda_{\mathrm{late}} \sum_t \ell_t
\;+\;
\lambda_{\mathrm{skip}} N_{\mathrm{skip}},
\]
where $N_{\mathrm{skip}}$ is the number of customer requests that remain unserved when the episode terminates. Thus, the metric jointly captures route efficiency, time-window violation, and service completeness. Throughout the experiments, we use $\lambda_{\mathrm{late}} = 1$ and $\lambda_{\mathrm{skip}} = 2$.

The term \emph{normalized} refers to the representation of the benchmark instances prior to inference. Customer coordinates are rescaled to $[0,1]$, demands are normalized by vehicle capacity, and all temporal quantities, including ready times, closing times, service durations, and appearance times, are normalized by the planning horizon; vehicle speed is rescaled consistently. As a result, the reported objective values are dimensionless normalized costs rather than raw physical distances or times.

\paragraph{Benchmark instances}
We employ two groups of benchmark instances. The first group comprises 176 static instances derived from Solomon-style topologies, organized by size and spatial distribution: 56 instances with 100 customers, 60 with 200 customers, and 60 with 400 customers. Each size level includes three topology types --- Random (R), Cluster (C), and a mixed Random-Cluster (RC) --- which impose different routing demands. Random instances test spatial dispersion, Cluster instances test local packing efficiency, and RC instances combine both characteristics. We will call this group the solomon-style benchmark for brevity, and refer to the individual datasets as \texttt{h100}, \texttt{h200}, and \texttt{h400}.

The second group consists of five synthetically generated datasets at increasing scales: \texttt{n20m1} (20 customers, 1 vehicle), \texttt{n50m3} (50 customers, 3 vehicles), \texttt{n100m5}, \texttt{n200m10}, and \texttt{n400m20}. These datasets scale both the number of customers and the fleet size simultaneously, enabling evaluation of coordination requirements under growing problem complexity. Each dataset contains 10,240 instances generated from the same distribution used during training. Each dataset is created by combining four degrees of dynamism (DoD = 0.10, 0.25, 0.50, 0.75) with four time-window ratios (0.25, 0.50, 0.75, 1.00). This group will be referred to as the PYTH benchmark, and the individual datasets will be referred to by their \texttt{nXmY} notation.

\paragraph{Baselines}
We compare COAST against several established methods. \textbf{MARDAM} is a multi-head attention-based routing policy that serves as a direct neural baseline with comparable capacity. \textbf{Attention Model (AM)} \cite{kool2018attention} and \textbf{PolyNet} \cite{hottung2025polynet} are Transformer-based routing architectures originally developed for static VRP, adapted here for the dynamic setting. We additionally include classical hyper-heuristic baselines: \textbf{GPHH} (genetic programming hyper-heuristic) \cite{TAM2026114234}, \textbf{C+C} (Clark-Wright savings with competence), \textbf{C+W} (Clark-Wright with waiting), and \textbf{WIQ+C} (work-in-queue plus competence). All neural baselines use the same training data distribution.

\paragraph{Training configuration}
Model-specific architectural settings and parameter counts are summarized in \autoref{tab:model_config}. The key point for interpretation is that COAST and all internal ablations share the same backbone, so the ablation study isolates the contribution of the proposed modules rather than capacity differences. By contrast, MARDAM, AM, and PolyNet should be viewed as external neural baselines with comparable or stronger capacity, making improvements over them conservative rather than attributable to an oversized backbone.

Training uses 500 epochs (1,000 iterations each) with a batch size of 512, so the number of training instances is 512,000. The learning rate is $10^{-4}$ via the Adam optimizer. The policy gradient objective employs a learned critic baseline with advantage normalization and entropy regularization ($\beta = 0.01$). All models are trained on DVRPTW instances with 50 customers and 3 vehicles; time windows are drawn uniformly from $[30, 91]$, demands from $[5, 41]$, and the degree of dynamism is sampled uniformly from $\{0.10, 0.25, 0.50, 0.75\}$ per instance, exposing all policies to low-, medium-, and high-dynamism regimes during training rather than tuning them to a single operating point.

\begin{table}[t]
\centering
\caption{Neural model configurations used in the experiments. The first row lists COAST and its internal ablations, which share the same backbone architecture. The subsequent rows list external neural baselines with comparable or stronger capacity.}
\label{tab:model_config}
\scriptsize
\resizebox{\linewidth}{!}{
\begin{tabular}{lcccccl}
\toprule
Model & Layers & Heads & FF & Dim & Params \\
\midrule
COAST + internal ablations & 2 & 4 & 256 & 128 & 0.649M  \\
MARDAM & 3 & 8 & 512 & 128 & 0.727M  \\
AM & 5 & 8 & 256 & 128 & 0.794M  \\
PolyNet & 5 & 8 & 256 & 128 & 0.830M  \\
\bottomrule
\end{tabular}}
\vspace{1mm}
\end{table}

\paragraph{Training dynamics}
\autoref{fig:learning_curves} shows the rolling-average validation cost over 500 training epochs for COAST, MARDAM, and representative ablation variants. COAST converges steadily and reaches its best validation cost by the final epoch, while MARDAM plateaus earlier at a consistently higher cost. Among the ablations, B5 (linear fusion) exhibits noticeably less stable convergence on larger instances, consistent with the degradation observed in the final performance results. All other ablation variants follow a convergence profile similar to COAST, indicating that the backbone is robust to component removal during optimization.

\begin{figure}[t]
\centering
\includegraphics[width=\linewidth]{figures/fig_learning_curves.pdf}
\caption{Validation cost curves during training. Lines show rolling-average validation cost per epoch for COAST, MARDAM, and selected ablation variants trained on \texttt{n50m3} instances.}
\label{fig:learning_curves}
\end{figure}



\subsection{Main results on Solomon-style datasets}

\autoref{tab:solomon_benchmark} reports the Solomon-style benchmark results grouped by instance size and topology. Instead of using only size-level averages, we separately report Random (R), Cluster (C), and mixed Random-Cluster (RC) instances because these topologies impose different routing challenges. Random instances emphasize global assignment and long-range coordination, Cluster instances favor compact intra-cluster routing, while RC instances require both local exploitation and global coordination.

At the smallest scale, \texttt{h100}, COAST(g) achieves the best mean cost on the Random and RC subsets, with costs of 51.33 and 65.69, respectively. The only exception is the Cluster subset, where MARDAM(g) slightly outperforms COAST(g), obtaining 65.08 compared with 67.65. This suggests that, for small and spatially concentrated instances, the benefit of the proposed decomposition is less pronounced because the routing structure is already relatively compact. Nevertheless, after aggregating all \texttt{h100} topologies, COAST(g) still achieves the best overall mean cost of 60.39.

As the instance size increases, the advantage of COAST becomes more consistent. On \texttt{h200}, COAST(g) obtains the lowest mean cost across all three topology groups, with an aggregate cost of 108.70 compared with 149.27 for MARDAM(g), 124.88 for AM(g), and 129.44 for PolyNet(g). The same trend is further amplified on \texttt{h400}, where COAST(g) again ranks first on Random, Cluster, and RC instances, leading to the best aggregate cost of 208.20. These results indicate that COAST scales more effectively as the number of customers grows, where route decomposition and coordination become increasingly important.

The topology-level results further show that the gains are not concentrated in a single spatial regime. At \texttt{h200}, MARDAM(g) is 37.4\%, 38.1\%, and 36.4\% higher than COAST(g) on the Random, Cluster, and RC groups, respectively, when measured relative to COAST(g). At \texttt{h400}, these gaps increase to 52.5\%, 46.8\%, and 58.8\%. The largest gap appears on the \texttt{h400}-RC subset, suggesting that COAST is particularly effective when both clustered local structure and dispersed global structure must be handled simultaneously.

The instance-level comparison in \autoref{tab:solomon_win_counts} confirms that the aggregate improvements are not caused by a small number of outlier instances. Across all 176 Solomon-style instances, COAST(g) outperforms MARDAM(g) on 157 instances, AM(g) on 166 instances, PolyNet(g) on 161 instances, and GPHH on 172 instances. COAST(g) also wins against C+C, C+W, and WIQ+C on all evaluated instances. The win rates become especially strong at larger scales: on both \texttt{h200} and \texttt{h400}, COAST(g) wins against MARDAM(g) on 100.0\% of the instances.

\autoref{tab:solomon_improvement} reports the average percentage improvement of COAST(g) over each baseline. Over all instances, COAST(g) improves upon MARDAM(g), AM(g), PolyNet(g), and GPHH by 22.4\%, 13.8\%, 13.9\%, and 33.3\%, respectively. The improvement over MARDAM(g) increases from 2.0\% on \texttt{h100} to 28.6\% on \texttt{h200} and 35.3\% on \texttt{h400}, further supporting the scalability of COAST. Similar trends are observed against the classical and handcrafted baselines, where the improvements exceed 60\% on average for C+C, C+W, and WIQ+C.

Comparing decoding strategies in \autoref{tab:solomon_benchmark}, COAST achieves strong results with both sampling and greedy decoding. The difference between COAST(s) and COAST(g) is generally small, and greedy decoding is often competitive or better in the aggregate rows. This indicates that COAST does not rely heavily on stochastic sampling to obtain high-quality solutions, which is desirable for efficient inference.

Overall, the results in \autoref{tab:solomon_benchmark}, \autoref{tab:solomon_win_counts}, and \autoref{tab:solomon_improvement} show that COAST is competitive at small scale and becomes increasingly advantageous on larger Solomon-style instances. Except for the \texttt{h100} Cluster subset, COAST consistently delivers the best greedy performance across topology groups. The high instance-level win rates and substantial percentage improvements further demonstrate that the proposed method provides robust and scalable gains across different spatial distributions. \autoref{fig:behavioral_analysis} provides additional evidence for the source of these gains by breaking down the contribution of each architectural component and illustrating how frequently the lookahead signal actively changes the dispatching decision.


\begin{table*}[t]
\centering
\caption{Solomon-style benchmark results grouped by dataset size and topology. Values are mean normalized costs over the evaluated scenarios in each row. The \texttt{All} row aggregates all topologies within the same dataset size, and the final row aggregates all 176 Solomon-style instances.}
\label{tab:solomon_benchmark}
\scriptsize
\resizebox{\textwidth}{!}{
\begin{tabular}{llrrrrrrrrrrrr}
\toprule
Dataset & Topology & COAST(s) & COAST(g) & MARDAM(s) & MARDAM(g) & AM(g) & AM(s) & PolyNet(g) & PolyNet(s) & GPHH & C+C & C+W & WIQ+C \\
\midrule
\texttt{h100} & Random & 52.73 & \textbf{51.33} & 56.53 & 54.89 & 61.22 & 60.28 & 59.57 & 98.89 & 66.39 & 144.02 & 144.40 & 178.61 \\
\texttt{h100} & Cluster & 66.68 & 67.65 & 66.11 & \textbf{65.08} & 73.01 & 72.56 & 70.36 & 100.38 & 80.63 & 130.71 & 118.12 & 175.02 \\
\texttt{h100} & RC & 66.57 & \textbf{65.69} & 68.90 & 67.42 & 73.07 & 71.80 & 70.31 & 112.22 & 80.11 & 152.99 & 141.10 & 181.66 \\
\texttt{h100} & All & 60.92 & \textbf{60.39} & 62.97 & 61.56 & 68.18 & 67.30 & 65.91 & 103.15 & 74.63 & 142.54 & 135.48 & 178.39 \\
\midrule
\texttt{h200} & Random & 106.03 & \textbf{103.29} & 145.03 & 141.90 & 122.33 & 121.73 & 120.81 & 169.48 & 188.74 & 270.39 & 262.94 & 346.91 \\
\texttt{h200} & Cluster & 118.06 & \textbf{117.61} & 163.22 & 162.43 & 126.64 & 127.77 & 137.04 & 171.67 & 191.74 & 269.58 & 248.67 & 361.24 \\
\texttt{h200} & RC & \textbf{104.97} & 105.19 & 146.90 & 143.48 & 125.66 & 126.27 & 130.48 & 183.64 & 181.00 & 282.86 & 283.12 & 350.99 \\
\texttt{h200} & All & 109.69 & \textbf{108.70} & 151.72 & 149.27 & 124.88 & 125.26 & 129.44 & 174.93 & 187.16 & 274.28 & 264.91 & 353.05 \\
\midrule
\texttt{h400} & Random & 208.86 & \textbf{200.71} & 317.68 & 306.15 & 237.63 & 238.96 & 232.22 & 349.95 & 356.53 & 613.33 & 600.64 & 716.09 \\
\texttt{h400} & Cluster & 225.12 & \textbf{220.65} & 339.35 & 323.84 & 248.35 & 251.41 & 258.78 & 361.37 & 374.85 & 630.06 & 616.13 & 752.28 \\
\texttt{h400} & RC & 208.20 & \textbf{203.26} & 341.53 & 322.73 & 250.32 & 253.20 & 244.98 & 405.95 & 352.72 & 634.84 & 636.76 & 724.18 \\
\texttt{h400} & All & 214.06 & \textbf{208.20} & 332.86 & 317.57 & 245.43 & 247.86 & 245.33 & 372.42 & 361.37 & 626.07 & 617.85 & 730.85 \\
\midrule
All & All & 129.75 & \textbf{127.25} & 185.23 & 178.74 & 147.94 & 148.61 & 148.73 & 219.42 & 210.74 & 352.29 & 344.05 & 426.27 \\
\bottomrule
\end{tabular}}
\end{table*}


\begin{figure*}[t]
\centering
\includegraphics[width=0.85\textwidth]{figures/fig_behavioral_analysis.pdf}
\caption{Behavioral analysis of COAST. Left: proportion of timesteps where the lookahead signal overrides the raw attention ranking, broken down by evaluation regime. Right: per-component cost contribution relative to the B0 baseline (no memory, ownership, or lookahead), showing the marginal gain from adding each proposed module.}
\label{fig:behavioral_analysis}
\end{figure*}


\begin{table*}[t]
\centering
\caption{Instance-level win counts and win rates of COAST(g) on the Solomon-style benchmark. Each entry has the format wins/instances (win rate), where a win indicates that COAST(g) obtains a lower normalized cost than the corresponding baseline.}
\label{tab:solomon_win_counts}
\scriptsize
\resizebox{\textwidth}{!}{
\begin{tabular}{lrrrrrrr}
\toprule
Dataset & MARDAM(g) & AM(g) & PolyNet(g) & GPHH & C+C & C+W & WIQ+C \\
\midrule
\texttt{h100} & 37/56 (66.1\%) & 48/56 (85.7\%) & 43/56 (76.8\%) & 53/56 (94.6\%) & 56/56 (100.0\%) & 56/56 (100.0\%) & 56/56 (100.0\%) \\
\texttt{h200} & 60/60 (100.0\%) & 58/60 (96.7\%) & 59/60 (98.3\%) & 59/60 (98.3\%) & 60/60 (100.0\%) & 60/60 (100.0\%) & 60/60 (100.0\%) \\
\texttt{h400} & 60/60 (100.0\%) & 60/60 (100.0\%) & 59/60 (98.3\%) & 60/60 (100.0\%) & 60/60 (100.0\%) & 60/60 (100.0\%) & 60/60 (100.0\%) \\
All & 157/176 (89.2\%) & 166/176 (94.3\%) & 161/176 (91.5\%) & 172/176 (97.7\%) & 176/176 (100.0\%) & 176/176 (100.0\%) & 176/176 (100.0\%) \\
\bottomrule
\end{tabular}}
\end{table*}


\begin{table}[t]
\centering
\caption{Average percentage improvement of COAST(g) over each baseline on the Solomon-style benchmark. Values are computed per instance as $(\mathrm{baseline}-\mathrm{COAST})/\mathrm{baseline}\times 100$ and then averaged within each group.}
\label{tab:solomon_improvement}
\scriptsize
\setlength{\tabcolsep}{3.2pt}
\resizebox{\columnwidth}{!}{
\begin{tabular}{lrrrrrrr}
\toprule
Dataset & MARDAM(g) & AM(g) & PolyNet(g) & GPHH & C+C & C+W & WIQ+C \\
\midrule
\texttt{h100} & 2.0 & 11.0 & 8.0 & 18.3 & 56.9 & 53.1 & 66.1 \\
\texttt{h200} & 28.6 & 14.0 & 17.3 & 40.0 & 61.1 & 59.5 & 69.4 \\
\texttt{h400} & 35.3 & 16.3 & 16.2 & 40.6 & 67.4 & 66.9 & 71.7 \\
All & 22.4 & 13.8 & 13.9 & 33.3 & 61.9 & 60.0 & 69.1 \\
\bottomrule
\end{tabular}}
\end{table}


\subsection{Main results on PYTH datasets}

\autoref{tab:pyth_benchmark} reports the results on the five synthetically generated PYTH datasets. Complementing the Solomon-style experiments, which evaluate robustness across spatial topologies, the PYTH benchmark focuses on scalability by increasing both the number of customers and the fleet size. This setting is therefore useful for examining whether the advantage of COAST persists when route construction and inter-vehicle coordination become increasingly difficult.

Across all PYTH scales, COAST(g) achieves the lowest normalized cost among the neural baselines. On the smallest single-vehicle setting, \texttt{n20m1}, COAST(g) obtains 24.78, slightly improving over MARDAM(g), AM(g), and PolyNet(g). Since this setting contains only one vehicle, it provides limited evidence for fleet-level coordination and mainly tests basic route construction. The comparison becomes more informative from \texttt{n50m3} onward, where multiple vehicles must be coordinated. In these multi-vehicle settings, COAST(g) consistently outperforms the baselines, achieving costs of 43.92, 94.85, 182.38, and 355.99 on \texttt{n50m3}, \texttt{n100m5}, \texttt{n200m10}, and \texttt{n400m20}, respectively.

The performance gap generally widens as the problem scale increases. Relative to COAST(g), MARDAM(g) is 5.4\%, 5.8\%, 4.9\%, 11.4\%, and 20.9\% higher from \texttt{n20m1} to \texttt{n400m20}, respectively. A similar advantage is observed over AM(g) and PolyNet(g), whose costs are consistently higher across all scales; on \texttt{n50m3}, for example, AM(g) and PolyNet(g) remain 12.9\% and 13.7\% above COAST(g). This trend supports the same scalability conclusion observed in the Solomon-style benchmark: COAST remains competitive on smaller instances, but its advantage becomes clearer when the routing problem requires stronger coordination over a larger set of customers and vehicles.

Comparing decoding strategies, COAST(g) is competitive with or better than COAST(s) on four of the five scales and is already stronger from \texttt{n50m3} onward. Sampling is only marginally better on \texttt{n20m1}, where the problem reduces to a compact single-route regime with little coordination burden. This again suggests that COAST does not depend heavily on stochastic sampling to obtain high-quality solutions, making greedy decoding an efficient and reliable inference option for the larger multi-vehicle settings that matter most in this study.

Overall, the PYTH results reinforce the main finding from the Solomon-style experiments. COAST provides stable gains across different benchmark settings, and its advantage becomes more pronounced as routing complexity increases. While the Solomon-style results demonstrate robustness across spatial distributions, the PYTH results show that COAST also scales effectively when both customer count and fleet size grow.

\begin{table*}[t]
\centering
\caption{PYTH benchmark results. Values are mean normalized cost $\pm$ standard deviation.}
\label{tab:pyth_benchmark}
\scriptsize
\resizebox{\textwidth}{!}{
\begin{tabular}{lrrrrrrrr}
\toprule
Dataset & COAST(s) & COAST(g) & MARDAM(s) & MARDAM(g) & AM(s) & AM(g) & PolyNet(s) & PolyNet(g) \\
\midrule
\texttt{n20m1} & $\textbf{24.72}\pm2.40$ & $24.78\pm2.38$ & $26.10\pm3.09$ & $26.13\pm3.16$ & $25.74\pm3.13$ & $25.40\pm2.75$ & $31.84\pm4.08$ & $25.58\pm2.94$ \\
\texttt{n50m3} & $44.10\pm4.09$ & $\textbf{43.92}\pm4.09$ & $46.85\pm4.71$ & $46.48\pm4.76$ & $50.28\pm5.49$ & $49.58\pm5.00$ & $68.57\pm8.22$ & $49.93\pm5.31$ \\
\texttt{n100m5} & $95.49\pm5.88$ & $\textbf{94.85}\pm5.89$ & $101.06\pm6.70$ & $99.50\pm6.56$ & $106.18\pm8.06$ & $103.85\pm7.14$ & $144.53\pm12.87$ & $105.13\pm8.42$ \\
\texttt{n200m10} & $184.70\pm8.95$ & $\textbf{182.38}\pm8.90$ & $209.02\pm14.70$ & $203.24\pm14.43$ & $201.54\pm13.02$ & $196.53\pm10.92$ & $293.54\pm24.83$ & $200.89\pm14.36$ \\
\texttt{n400m20} & $363.35\pm13.82$ & $\textbf{355.99}\pm13.43$ & $443.86\pm34.78$ & $430.48\pm35.05$ & $389.99\pm21.35$ & $378.98\pm17.10$ & $558.21\pm41.05$ & $389.72\pm26.42$ \\
\bottomrule
\end{tabular}}
\end{table*}

\subsection{Effect of dynamism}

To isolate the effect of online uncertainty, we evaluate a dedicated dynamic benchmark grid built from the same \texttt{n50m3} regime used in training: 50 customers, 3 vehicles, horizon 480, and greedy decoding. The grid contains 512 instances organized as 16 cells, obtained by combining four degrees of dynamism (DoD = 0.10, 0.25, 0.50, 0.75) with four time-window ratios (0.25, 0.50, 0.75, 1.00), with 32 instances per cell. \autoref{tab:dynamic_summary} reports the aggregated cell averages by DoD and by time-window ratio, while \autoref{fig:dynamic_sensitivity} provides cell-level heatmaps over the full DoD $\times$ TW grid for all four neural methods.

The full grid confirms that COAST is the best-performing method in all 16 DoD--TW cells. Averaged over the entire grid, COAST attains a mean normalized cost of 43.80, compared with 45.02 for MARDAM, 47.15 for PolyNet, and 47.65 for AM. The best COAST cell occurs at DoD = 0.50 and TW = 0.25 with cost 41.21, while the most difficult cell is DoD = 0.75 and TW = 1.00 with cost 46.88. Even in that hardest setting, COAST remains better than MARDAM (47.96), PolyNet (49.44), and AM (49.53), indicating that the advantage is not confined to easy operating points.

As DoD increases from 0.10 to 0.75, COAST remains comparatively stable, with mean cost increasing only from 43.64 to 44.45, a range of 0.81. Over the same sweep, MARDAM increases from 44.07 to 46.68, a much larger range of 2.61, while PolyNet and AM average 47.78 to 46.85 and 48.08 to 47.29, respectively. COAST therefore preserves the best absolute performance throughout the entire uncertainty range, and its advantage over MARDAM widens from 1.0\% at DoD = 0.10 to 5.0\% at DoD = 0.75 when averaged over time-window settings. At the cell level, the gap to MARDAM ranges from only 0.5\% at (DoD = 0.10, TW = 0.50) to 6.7\% at (DoD = 0.75, TW = 0.25), which is consistent with the intended role of the ownership and lookahead signals: when more customers are revealed online, the policy benefits more from explicit coordination and future-aware scoring.

All methods also degrade as the time-window ratio increases. Averaged over DoD, COAST rises from 41.71 at TW = 0.25 to 45.90 at TW = 1.00, while MARDAM rises from 43.51 to 46.62, PolyNet from 44.01 to 50.76, and AM from 44.83 to 50.73. This confirms that tighter windows reduce the feasible action set for every model, but the degradation is consistently milder for COAST than for the stronger neural baselines. The gap to AM, for example, is 7.5\% at TW = 0.25 and remains 10.5\% at TW = 1.00; against PolyNet, the corresponding gaps are 5.5\% and 10.6\%. Taken together, these results show that the proposed architecture is not only more robust to increasing dynamism, but also better at preserving route quality when temporal feasibility becomes more restrictive.

\begin{table}[t]
\centering
\caption{Dynamic benchmark summary aggregated by degree of dynamism and time-window ratio. Values are mean normalized cost over the corresponding four cells.}
\label{tab:dynamic_summary}
\scriptsize
\setlength{\tabcolsep}{4pt}
\resizebox{\columnwidth}{!}{
\begin{tabular}{llrrrr}
\toprule
Factor & Level & COAST & MARDAM & PolyNet & AM \\
\midrule
DoD & 0.10 & \textbf{43.64} & 44.07 & 47.78 & 48.08 \\
DoD & 0.25 & \textbf{43.50} & 44.35 & 47.25 & 47.68 \\
DoD & 0.50 & \textbf{43.61} & 44.99 & 46.73 & 47.55 \\
DoD & 0.75 & \textbf{44.45} & 46.68 & 46.85 & 47.29 \\
\midrule
TW & 0.25 & \textbf{41.71} & 43.51 & 44.01 & 44.83 \\
TW & 0.50 & \textbf{43.15} & 44.37 & 45.89 & 46.73 \\
TW & 0.75 & \textbf{44.43} & 45.59 & 47.94 & 48.31 \\
TW & 1.00 & \textbf{45.90} & 46.62 & 50.76 & 50.73 \\
\bottomrule
\end{tabular}}
\end{table}

\begin{figure}[t]
\centering
\includegraphics[width=\linewidth]{figures/fig_dynamic_sensitivity.pdf}
\caption{Cell-level heatmaps on the dynamic benchmark grid. Each panel shows the mean normalized cost of one neural method for every DoD $\times$ TW configuration; lower values are better (lighter is better).}
\label{fig:dynamic_sensitivity}
\end{figure}

\subsection{Ablation study and hypothesis analysis}

The proposed method is built around four architectural claims: \textbf{(H1)} explicit ownership improves fleet coordination, \textbf{(H2)} candidate-conditioned lookahead reduces myopic dispatching, \textbf{(H3)} edge-aware scoring improves constraint handling, and \textbf{(H4)} non-linear fusion is needed to combine these signals robustly. To test these claims, we design ablation variants that each isolate one component. \textbf{B0} removes all structured signals (no memory, ownership, or lookahead), leaving only the base attention mechanism; \textbf{B1} retains memory-based vehicle state but removes ownership and lookahead; \textbf{B3} includes only the lookahead signal without coordination; \textbf{B5} retains all signals but replaces the learned non-linear fusion with a linear combination; \textbf{EdgeOff} disables explicit edge features while keeping all other components; and two clean ablations remove either the ownership head or the lookahead head individually. Each hypothesis maps directly to one or more of these variants, so the ablation results and hypothesis tests are presented together.

The full COAST model demonstrates consistent advantages once the problem involves sufficient customers and vehicles for coordination to matter. \autoref{tab:ablation_pyth} reports results for all variants on the five PYTH scales; \autoref{fig:ablation_waterfall} visualizes the cumulative cost reduction from each component relative to B0. On the smallest \texttt{n20m1} configuration, several simplified variants perform competitively, confirming that this regime lacks the multi-vehicle dynamics necessary to distinguish architectural choices. From \texttt{n50m3} to \texttt{n400m20}, COAST outperforms all ablated variants, and the margin widens with scale.

\begin{table*}[t]
\centering
\caption{Ablation results on PYTH datasets. Values are mean normalized cost $\pm$ standard deviation.}
\label{tab:ablation_pyth}
\scriptsize
\resizebox{\textwidth}{!}{
\begin{tabular}{lrrrrrrrr}
\toprule
Dataset & COAST & B0 & B1 & B3 & B5 & EdgeOff & no\_ownership & no\_lookahead \\
\midrule
\texttt{n20m1} & $24.78\pm2.38$ & $\textbf{23.51}\pm2.36$ & $23.81\pm2.43$ & $23.45\pm2.43$ & $23.79\pm2.42$ & $24.91\pm2.34$ & $23.37\pm2.46$ & $23.96\pm2.49$ \\
\texttt{n50m3} & $\textbf{43.92}\pm4.09$ & $45.28\pm4.26$ & $45.43\pm4.32$ & $44.91\pm4.25$ & $45.49\pm4.31$ & $45.43\pm4.20$ & $45.20\pm4.25$ & $44.74\pm4.22$ \\
\texttt{n100m5} & $\textbf{94.85}\pm5.89$ & $97.35\pm6.17$ & $97.71\pm6.36$ & $96.86\pm6.15$ & $101.67\pm9.02$ & $97.25\pm5.99$ & $97.44\pm6.27$ & $96.44\pm6.17$ \\
\texttt{n200m10} & $\textbf{182.38}\pm8.90$ & $186.52\pm9.21$ & $186.26\pm9.63$ & $184.75\pm9.13$ & $259.10\pm34.30$ & $186.94\pm8.90$ & $186.22\pm9.26$ & $184.33\pm9.43$ \\
\texttt{n400m20} & $\textbf{355.99}\pm13.43$ & $363.43\pm13.86$ & $361.33\pm14.61$ & $357.81\pm13.77$ & $515.72\pm67.45$ & $365.43\pm13.16$ & $362.11\pm13.56$ & $357.18\pm14.44$ \\
\bottomrule
\end{tabular}}
\end{table*}

For H1 and H2, the cleanest tests are direct removals of the ownership and lookahead modules, respectively. For H3, the relevant comparison is COAST versus EdgeOff, especially under the tighter time-window and higher-dynamism cells where feasibility pressure is strongest. For H4, the key comparison is COAST versus B5, which keeps all signals but replaces the learned MLP fusion with a linear combination. We draw evidence from \autoref{tab:ablation_pyth}, the dynamic-grid analysis in \autoref{tab:dynamic_summary} and \autoref{fig:dynamic_sensitivity}, and the complementary view in \autoref{fig:ablation_waterfall}. Throughout the discussion below, we emphasize \texttt{n50m3} and larger instances, because \texttt{n20m1} contains only one vehicle and cannot serve as strong evidence for coordination-oriented claims.

\textbf{H1 (ownership improves coordination)} asks whether the explicit ownership mechanism contributes beyond the shared policy backbone. The evidence is strongest in the clean COAST versus no-ownership comparison: from \texttt{n50m3} onward, removing ownership increases normalized cost by 1.7--2.9\%. This confirms that the per-vehicle coordination memory and learned ownership bias contribute to more efficient fleet allocation once multiple vehicles must compete for dynamically revealed customers.

\textbf{H2 (lookahead reduces myopia)} tests whether scoring candidate actions with a future-impact signal improves decision quality. This hypothesis receives positive but more modest support. Removing lookahead consistently increases cost on medium and large PYTH datasets, but the effect remains smaller (0.3--1.9\%) than those observed for ownership removal or fusion simplification. The directional trend is therefore stable, but the relatively small magnitude suggests that lookahead helps as a refinement signal rather than as the dominant driver of performance.

\textbf{H3 (edge features improve feasibility)} evaluates whether exposing route-level and constraint-level information directly to the scorer is beneficial. The EdgeOff ablation is consistently worse than COAST, with a 2.5--3.4\% increase on the PYTH scale experiments, and the dynamic benchmark further shows that COAST preserves a clearer advantage under tighter time windows and higher dynamism. Taken together, these results support the claim that explicit edge-level features improve the policy's ability to act under feasibility pressure rather than relying only on latent node embeddings.

\textbf{H4 (non-linear fusion improves robustness)} concerns the interaction among the three signals rather than any signal in isolation. This is the strongest hypothesis in the paper. The B5 ablation, which preserves the same signal sources but replaces the learned non-linear fusion with a linear combination, incurs only modest degradation on very small instances but collapses on larger ones (42.1\% increase on \texttt{n200m10} and 44.9\% on \texttt{n400m20}). This indicates that the signals are not merely additive: their relative importance depends on context, and a linear rule cannot express the required trade-offs between immediate compatibility, ownership, and downstream routing consequences.

\begin{figure}[t]
\centering
\includegraphics[width=\linewidth]{figures/fig_ablation_waterfall.pdf}
\caption{Cost increase of each ablation variant relative to COAST (\%) across four PYTH scales under greedy decoding. B5 (linear fusion) collapses at large scale (42--45\% on \texttt{n200m10}/\texttt{n400m20}), while all other ablations remain consistently within 1--4\%.}
\label{fig:ablation_waterfall}
\end{figure}

\autoref{tab:hypotheses} summarizes the mapping from each hypothesis to its comparison, evidence, and conclusion.

\begin{table*}[t]
\centering
\caption{Hypothesis evaluation summary. Positive gaps indicate that the ablated variant incurs higher cost than COAST.}
\label{tab:hypotheses}
\footnotesize
\resizebox{\textwidth}{!}{
\begin{tabular}{llll}
\toprule
Hypothesis & Comparison & Main evidence & Conclusion \\
\midrule
H1: ownership improves coordination & COAST vs no\_ownership & no\_ownership is 1.7--2.9\% worse on \texttt{n50m3}--\texttt{n400m20} & Supported on medium/large instances \\
H2: lookahead reduces myopia & COAST vs no\_lookahead & no\_lookahead is 0.3--1.9\% worse on \texttt{n50m3}--\texttt{n400m20} & Mildly supported; effect is smaller \\
H3: edge features improve feasibility & COAST vs EdgeOff & EdgeOff is 2.5--3.4\% worse on PYTH scale; tight-TW cost is 48.00 vs 46.13 & Supported, especially under tight constraints \\
H4: nonlinear fusion improves robustness & COAST vs B5 & B5 is 7.2\% worse on \texttt{n100m5}, 42.1\% on \texttt{n200m10}, and 44.9\% on \texttt{n400m20} & Strongly supported at scale \\
\bottomrule
\end{tabular}}
\end{table*}

Overall, the hypothesis tests support the central decomposition strategy, but with different strengths. H4 is the clearest architectural result because the degradation is both large and strongly scale-dependent. H1 and H3 provide consistent medium-strength evidence across the regimes where coordination and feasibility are genuinely challenging. H2 is directionally supported, but the weaker margin suggests that improving the training signal for lookahead remains an important direction for future work.

Overall, the results demonstrate that COAST is most effective in the regimes targeted by its architecture: medium-to-large dynamic routing problems where vehicle coordination, temporal feasibility, and action-conditioned anticipation interact. The small single-vehicle regime (\texttt{n20m1}) is less favorable to the full model, which is expected given that coordination requirements are minimal in this setting.

\subsection{Limitations}

Several limitations of the current study should be acknowledged. First, the reported results are based on single-seed training for each model configuration. The smaller effect sizes observed for the lookahead and ownership ablations could be sensitive to training variance; multi-seed experiments with statistical significance testing would strengthen confidence in these findings. Second, the lookahead head is trained solely through the policy gradient objective; a supervised calibration target could potentially enhance its effectiveness and provide a clearer signal for validating the anticipation hypothesis. Third, the current evaluation focuses on normalized routing cost as the primary metric; a more detailed decomposition into travel distance, skipped requests, and time-window penalties would provide additional insight into behavioral differences between methods. Finally, classical optimization baselines such as LKH3 and OR-Tools, while informative for static or oracle settings, are not directly applicable to the dynamic sequential decision framework without substantial modification; their inclusion remains an avenue for future work.

\section{Conclusions and Future works} \label{sec:futurework}
\bibliographystyle{plain}
\bibliography{ref}
\printglossaries
\end{document}
